\ChapterImageStar[cap:resultados]{Resultados}{./images/fondo.png}\label{cap:resultados}
\mbox{}\\

\section{Caracterización del Grupo \GRID\ e identificación de necesidades}
\noindent

La caracterización realizada sobre la infraestructura del grupo \GRID\ permitió identificar las capacidades tecnológicas existentes, los 
mecanismos actuales de protección y las necesidades asociadas al fortalecimiento de la seguridad en entornos de computación en la nube. Se 
evidenció que el grupo dispone de una infraestructura virtualizada soportada mediante hipervisores, servicios de red segmentados y
mecanismos básicos de protección perimetral, los cuales sirven como base para la integración de una arquitectura de defensa en profundidad.

El análisis de \textit{stakeholders} permitió identificar como actores principales al grupo \GRID, docentes y estudiantes del programa de 
Ingeniería de Sistemas y Computación, estableciendo su relación con los procesos de investigación, docencia y administración de los 
recursos tecnológicos. Asimismo, se determinó la necesidad de fortalecer los mecanismos de supervisión, monitoreo y gestión de eventos de 
seguridad dentro de la infraestructura, con el propósito de mejorar la trazabilidad y visibilidad operacional del entorno.

A partir de los hallazgos obtenidos en el estudio de Mapeo Sistemático de Literatura (\SMS) y del análisis \DAR\ realizado sobre diferentes 
tecnologías de seguridad, se seleccionaron herramientas orientadas a fortalecer las capacidades de detección, correlación y respuesta ante 
incidentes dentro de la infraestructura cloud del \GRID. Como resultado, se definió una arquitectura basada en la integración de un 
firewall \NGFW\ SonicWall NSA 6600, el \IDPS\ Suricata y el \SIEM\ Wazuh, estableciendo un modelo interoperable alineado con controles de las 
normas \ISO/IEC 27001 e \ISO/IEC 27017.

\section{Revisión sistemática de la literatura}
\noindent

El Estudio de Mapeo Sistemático de Literatura (\SMS) permitió identificar, analizar y clasificar un conjunto de 181 estudios primarios 
relacionados con arquitecturas de seguridad, controles y mecanismos de protección en entornos de computación en la nube. El proceso 
metodológico aplicado contempló una estrategia híbrida de búsqueda basada en consultas automáticas en bases de datos científicas y la 
técnica de \textit{snowballing}, garantizando trazabilidad, reproducibilidad y rigurosidad en la selección de los estudios analizados.

Los resultados obtenidos evidenciaron un crecimiento significativo en la producción científica relacionada con ciberseguridad en la nube 
durante el periodo comprendido entre 2022 y 2025, identificando un enfoque recurrente hacia arquitecturas Zero Trust, mecanismos de 
monitoreo continuo, correlación de eventos, protección perimetral y supervisión de infraestructuras distribuidas. Asimismo, los estudios 
seleccionados mostraron una amplia adopción de estándares y buenas prácticas alineadas con las normas \ISO/IEC 27001 e \ISO/IEC 27017.

De igual forma, el \SMS permitió establecer una estructura taxonómica orientada a clasificar los estudios según su tipo de contribución, 
enfoque de seguridad y contexto de aplicación, proporcionando una visión integral sobre las tendencias actuales en arquitecturas de 
seguridad para entornos cloud. Esta clasificación facilitó la identificación de patrones tecnológicos y metodológicos presentes en la 
literatura, así como la relación existente entre monitoreo, detección de amenazas, gestión de eventos y protección de activos de 
información dentro de infraestructuras institucionales y académicas.

\section{Análisis de Decisiones y Resolución (\DAR)}
\noindent

La aplicación de la metodología \DAR\ (\textit{Decision Analysis and Resolution}) del modelo \CMMI\ permitió evaluar de manera estructurada 
diferentes tecnologías orientadas a fortalecer la arquitectura de seguridad propuesta para el entorno \GRID. El proceso de evaluación se 
fundamentó en criterios técnicos, operativos y organizacionales alineados con las necesidades, problemas y oportunidades identificadas 
previamente, así como con controles establecidos en las normas \ISO/IEC 27001 e \ISO/IEC 27017.

Entre los criterios considerados se incluyeron aspectos como rentabilidad, tipo de licenciamiento, facilidad de implementación, calidad de 
la documentación, escalabilidad, consumo de recursos, capacidad de detección, integración con otras herramientas, actualización de reglas y 
capacidad de correlación de eventos. Estos criterios permitieron establecer un análisis comparativo orientado a determinar la viabilidad y 
pertinencia de cada alternativa tecnológica dentro del contexto de la infraestructura cloud del grupo \GRID.

Como resultado del proceso \DAR, las tecnologías \textbf{Suricata} y \textbf{Wazuh} obtuvieron las puntuaciones más altas dentro de las 
categorías \IDPS\ y \SIEM\ respectivamente, evidenciando una mayor compatibilidad con los requerimientos técnicos y operativos definidos 
para la arquitectura de defensa en profundidad propuesta. Los resultados obtenidos reflejan que ambas herramientas ofrecen capacidades 
robustas de monitoreo, detección, correlación de eventos y respuesta ante incidentes, además de una adecuada interoperabilidad con el 
ecosistema tecnológico planteado para el entorno \GRID.

\subsection{Arquitectura de defensa en profundidad}
\noindent

Se especificó una arquitectura de defensa en profundidad orientada al fortalecimiento de las capacidades de monitoreo, detección, 
correlación de eventos y respuesta ante incidentes dentro de la infraestructura cloud del grupo \GRID. La arquitectura integra componentes 
físicos, virtualizados y servicios de seguridad interoperables, conformando un modelo multicapa alineado con controles de las normas 
\ISO/IEC 27001 e \ISO/IEC 27017.

El diseño arquitectónico fue desarrollado mediante el marco de trabajo TOGAF y modelado utilizando ArchiMate como lenguaje de representación 
empresarial. Como resultado, se construyeron siete vistas arquitectónicas especializadas, orientadas a representar la interacción entre las 
capas de negocio, aplicación y tecnología, así como la cooperación entre actores, procesos, servicios, mecanismos de monitoreo y componentes 
de infraestructura involucrados en la operación del entorno \GRID.

La arquitectura incorpora mecanismos de segmentación de red, monitoreo continuo, inspección de tráfico, correlación de eventos y respuesta 
activa, integrando tecnologías como el firewall SonicWall NSA 6600, el \IDPS\ Suricata y el \SIEM\ Wazuh. Asimismo, se modeló el flujo 
operativo de la infraestructura mediante Draw.io, permitiendo representar de manera holística las zonas de red, el tráfico entre 
componentes y la trazabilidad operacional existente entre los diferentes elementos que conforman la arquitectura propuesta.

\section{Implementación del Prototipo de Arquitectura de defensa en profundidad}
\noindent

La materialización de la arquitectura propuesta se realizó mediante un \PMV o prototipo funcional queintegra todos los componentes
seleccionados y diseñados para la debida realización de esta sección.

\subsection{Implementación del Firewall Perimetral OPNsense}

La implementación del firewall OPNsense permitió establecer mecanismos de segmentación lógica, control de acceso y filtrado del tráfico 
dentro del laboratorio virtual. Mediante la configuración de interfaces virtualizadas independientes y reglas restrictivas de comunicación, 
se logró aislar los diferentes segmentos de red correspondientes a administración, monitoreo y servidores.

Asimismo, el despliegue de reglas de firewall sobre la interfaz WAN permitió establecer capacidades de control perimetral y soporte para 
procesos de respuesta activa provenientes de la herramienta \SIEM\ Wazuh. La implementación de aliases dinámicos permitió automatizar 
procesos de bloqueo temporal sobre direcciones IP identificadas como potencialmente maliciosas.

\subsection{Implementación del Sistema \IDPS\ Suricata}

La implementación del sistema \IDPS\ Suricata permitió realizar inspección profunda del tráfico de red mediante técnicas de port mirroring 
configuradas sobre la infraestructura virtualizada. Esto posibilitó analizar paquetes en tiempo real sin afectar el desempeño operativo de 
las máquinas desplegadas.

A partir de la utilización de firmas de detección y reglas de inspección, se logró identificar actividades sospechosas, tráfico 
potencialmente malicioso y eventos asociados a amenazas de red. Los registros generados en el archivo \textit{eve.json} fueron 
posteriormente centralizados por Wazuh para su correlación y monitoreo continuo.

\subsection{Implementación de la Plataforma \SIEM\ Wazuh}

La implementación de la plataforma \SIEM\ Wazuh permitió centralizar registros provenientes de los agentes desplegados sobre las diferentes 
máquinas del laboratorio, incluyendo el firewall OPNsense, el \IDPS\ Suricata y la máquina objetivo.

Mediante el uso de agentes y configuraciones personalizadas en el archivo \textit{ossec.conf}, se habilitaron capacidades de monitoreo 
continuo, clasificación de alertas, correlación de eventos y visualización centralizada de incidentes de seguridad.

Adicionalmente, se desarrollaron dashboards personalizados orientados al análisis de severidad de alertas, monitoreo de eventos por agente 
y visualización de firmas detectadas por Suricata, facilitando los procesos de auditoría y supervisión de la infraestructura.

\section{Validación técnica y funcional}
\noindent

Una vez implementados los componentes que conforman la arquitectura de seguridad propuesta en el prototipo, se procedió a 
ejecutar un conjunto de escenarios de validación orientados a evaluar su funcionamiento dentro del entorno prototipo. El 
propósito de esta fase fue verificar la capacidad de las tecnologías implementadas para detectar, registrar y gestionar 
eventos de seguridad, así como comprobar la interoperabilidad entre los diferentes componentes de la arquitectura.

Para ello, se desarrollaron pruebas controladas que simularon diferentes situaciones potencialmente relevantes para la 
seguridad de la infraestructura, permitiendo analizar el comportamiento de los mecanismos de protección perimetral, 
detección de intrusiones y gestión centralizada de eventos. Los resultados obtenidos permitieron validar la viabilidad 
técnica de la propuesta y evidenciar el cumplimiento de los objetivos definidos para la arquitectura de seguridad del \GRID.

\section{Métricas de impacto y capacidades resultantes}
\noindent

La implementación exitosa del prototipo permitió fortalecer las capacidades de seguridad y monitoreo del laboratorio virtual del grupo 
\GRID\ mediante la integración de tecnologías de código abierto orientadas a la protección de infraestructuras de computación en la nube.

\begin{itemize}
    \item \textbf{Centralización de eventos de seguridad:} Consolidación de registros provenientes del firewall, el \IDPS\ y las máquinas monitoreadas mediante la plataforma \SIEM\ Wazuh.

    \item \textbf{Capacidades de monitoreo continuo:} Implementación de dashboards personalizados para visualización de alertas, análisis de severidad y seguimiento de eventos de seguridad en tiempo real.

    \item \textbf{Inspección profunda de tráfico:} Despliegue del sistema \IDPS\ Suricata mediante técnicas de port mirroring para análisis de tráfico interno y externo sin afectar el rendimiento de la infraestructura.

    \item \textbf{Respuesta activa automatizada:} Integración entre Wazuh y OPNsense para aplicar bloqueos dinámicos sobre direcciones IP sospechosas detectadas por correlación de eventos.

    \item \textbf{Segmentación lógica de red:} Separación de entornos de administración, monitoreo y servicios mediante interfaces virtualizadas y políticas restrictivas de firewall.

    \item \textbf{Correlación de eventos de seguridad:} Relación automática entre alertas provenientes de diferentes fuentes para mejorar la detección de amenazas y el análisis de incidentes.

    \item \textbf{Cumplimiento de estándares internacionales:} Implementación alineada con controles de las normas \ISO/IEC 27001:2022, \ISO/IEC 27017:2015 y capacidades de monitoreo bajo \NISTCSF\ SP 800-53.

    \item \textbf{Escalabilidad del entorno:} Arquitectura virtualizada preparada para incorporar nuevas máquinas, agentes y mecanismos de seguridad sin afectar la operación del laboratorio.
\end{itemize}

\section{Contribuciones al conocimiento}
\noindent

Este proyecto genera múltiples contribuciones orientadas al fortalecimiento de la seguridad en infraestructuras de computación en la nube dentro de contextos académicos e institucionales:

\begin{itemize}
	\item \textbf{Metodológica}: Integración de un proceso reproducible compuesto por caracterización organizacional, mapeo sistemático de literatura, análisis \DAR\ y modelado arquitectónico basado en TOGAF y ArchiMate.
	
	\item \textbf{Técnica}: Diseño de una arquitectura de defensa en profundidad orientada a la supervisión, correlación y respuesta ante eventos de seguridad mediante tecnologías como SonicWall NSA 6600, Suricata y Wazuh.
	
	\item \textbf{Académica}: Consolidación y documentación de referentes metodológicos y tecnológicos relacionados con seguridad en la nube, proporcionando trazabilidad y soporte científico para futuras investigaciones.
	
	\item \textbf{Institucional}: Fortalecimiento de las capacidades de monitoreo, protección y gestión de eventos de seguridad dentro de la infraestructura del grupo \GRID, alineando la propuesta con controles de la \ISO/IEC 27001 e \ISO/IEC 27017.
\end{itemize}

Los resultados demuestran que la incorporación de los aspectos de seguridad seleccionados para la infraestructura del Grupo \GRID\ 
representa una solución viable y eficiente para la protección de los activos de información existentes, sin desestimar la inversión previa 
en tecnologías de seguridad y proporcionando una base sólida para el crecimiento futuro de la seguridad en entornos de computación en la 
nube institucionales.

% !TODO Validar este titulo con sepúlveda
%\section{Validación del proyecto}

%Con el objetivo de concluir formalmente el proyecto y verificar el cumplimiento de los objetivos planteados, se realizó una reunión entre 
%los integrantes del equipo. En dicho encuentro, y mediante consenso, se acordó dar por finalizado el proyecto.
